% Transcription — fonds Alexandre Grothendieck, shelfmark 117, batch 1 (pages 1-20).
% Established from the facsimile archives/batches/117/batch-01.pdf (a mirror of
% https://grothendieck.umontpellier.fr/117.pdf), per the skill
% transcribe-grothendieck. The batch file carries no cover sheet: PDF page N
% is page N of the folder (the Montpellier footer prints « 1 » ... « 20 »),
% and the archivists' pencil numbers bottom-left agree wherever legible
% (« 2 », « 3 », « 4 », « 5 », « 7 », « 9 », « 11 » ... « 20 »).
%
% Pass: Opus 5.5 (claude-opus-5-5), 2026-09-23 — first pass, unchecked against the pages by a human.
%
% Thirteen of the twenty pages are transcribed. Absent: pages 6, 8, 10, 12,
% 14, 16 and 18, the versos of the leaves carrying his pages 1-7 — ruled
% administrative tables in another hand (names, amounts, « Fixe / Points »),
% with nothing of his on any of them; absent under the settled rule for
% unrelated versos. Page 20 is a separate leaf and is transcribed.
%
% What the batch is.
%   1      A sheet of computer listing paper, landscape, headed in a box
%          « Structure des dérivateurs » and dated « déc. 83 » in the top
%          right corner (both his, both transcribed where they stand):
%          for a datum H(I), the three levels Hom ∈ Ens, Homhot ∈ Hot and
%          the internal Homhot ∈ Hot(I), then the adjunction f_! ⊣ f^* at
%          the three levels (1')-(3') and the dual f^* ⊣ f_* (1'')-(3''),
%          the internal forms boxed.
%   2      Scratch on the same listing paper: π_0, π_1 of an object of
%          H(Δ_1), 2-cells between two arrows x → y, J ⊂ I and H(I) → H(J),
%          « M W », « Hot », « Cl (M, W) ». Recorded compactly.
%   3-4    An unheaded programme in nine items a)-i) (letters of f)-i) overwritten): models, long exact
%          sequences, resolutions and derived functors, homotopy notions
%          on models, internal homotopies (Hom in Hot, « dans Quillen »),
%          operations and invariants (his margin: f) and g) « peu
%          compris »), K(H) via exact triangles, internal Hom and
%          multiplicative structure; then « Sauf peut-être ... f) et g),
%          il semble que tout puisse se décrire en termes de (M, W ⊂
%          Fl(M)). Dégager les bons axiomes ». Page 4's margin carries the
%          date « 24.11.83 », transcribed in its \marginal.
%   5, 7, 9, 11, 13, 15, 17, 19, 20
%          His own run, paginated 1 to 9 at the top centre (recorded once
%          per page in a \note): « Théorie (de variance) homotopique H ».
%          (1) H: A ↦ H(A), 2-contravariant, H(Δ_0) = Hot_H, the functor
%          α_A: H(A) → Hom(A, Hot_H), square (5); (Hot 1) H(∐A_i) ≃ ∏H(A_i);
%          (2) integration: a first version via L_A: H(A) → H(A^) and the
%          square (8), cancelled by diagonal strokes on page 7; then
%          (Hot 2) f^*_H has a left adjoint f^H_!, ∫_A = (π_A)_! plays the
%          role of lim→_A, ∫_A X = ∫_B f_!(X); (3) cointegration ∏_A =
%          (π_A)_*, (Hot 3) right adjoints f^H_*; duality H°(A) = H(A)°,
%          ∫ and ∏ exchanged; the additive case. Ex 1: H(A) = Hom(A, M) for
%          M complete and cocomplete. Ex 2: H(A) = W_A^{-1} Hom(A, M) for a
%          pair (M, W); a Remarque on f: (M, W) → (M', W') and when the
%          induced morphism commutes with f_!, f_* (an adjoint g with
%          f^{-1}(W') = W and fg(X') → X' ∈ W'); « théories spéciales »; the
%          question of comparing two (M_i, W_i) giving the same theory.
%          Ex 3: M = (Cat), W a « localisateur fondamental » — « J'ai prouvé
%          l'existence des f_!, mais non encore celle des f_* ». Ex 4:
%          complexes in an abelian category, the derived category of
%          sheaves on A°^ × T, Rf_* always, Lf_! at the level of sheaves of
%          sets. « Propriétés formelles »: (Hot 4) base change in a
%          cartesian square (f cofibrant / f fibrant, with his « ?? »),
%          (Hot 5) f fully faithful, (Hot 6) f a localisation. The list
%          breaks off at the foot of page 20 (his 9).
%
% Notation fixed for the batch. \mathcal{H} for his looped script H (the
% theory, and H(A)); \mathcal{X} for his looped X (an object of H(A));
% \mathrm{Hot}_{\mathcal H} for « Hot_H »; \mathrm{Hot} alone where he
% writes « Hot »; \underline{\mathrm{Hom}} for his underlined Hom (the
% functor category Hom(A, M) and, on page 1, the internal Homhot);
% \mathrm{Homhot} for « Hom hot », written as one word; f^{*}_{\mathcal H},
% f^{\mathcal H}_{!}, f^{\mathcal H}_{*} with the index placed as he places
% it (the star low for f^*, the H high for f_! and f_*); A^{\wedge} for his
% « A^ » (he also writes a hat, \widehat{A}, on page 7, kept); \Delta_0,
% \Delta_1 for his Δ with index; (\mathrm{Cat}) with his parentheses;
% \mathrm{Fl}(M) for « Fl(M) »; \varinjlim / \varprojlim for his lim with
% arrows. Underlined words are set in \emph. His « tt », « pl. fid. »,
% « évt. », « cat. », « resp. », « i.e. » are kept. His sign « = » for
% « même(s) » (as in folders 113 and 116) is expanded to « mêmes » with
% \add on page 17.
%
% Dates. The folder's \dating is the inventory's own, verbatim. The leaves
% themselves carry « déc. 83 » (page 1) and « 24.11.83 » (page 4, margin);
% both are transcribed where they stand, and nothing else in this file
% dates any page.
%
% Cross-references. Page 21 (batch 2) takes up again, as « Résultat des
% diagrammes », the square with (h_A)_! : H(A) → H(A^) that page 7 draws as
% (8) with L_A and then cancels.
\documentclass[11pt,a4paper]{article}
% Le préambule commun à toutes les transcriptions du fonds Grothendieck.
%
% Il tient en une idée : **l'incertitude se compose, elle ne se résout pas.**
% Une transcription qui lisse un mot douteux a détruit la seule chose pour
% laquelle on la faisait — un lecteur doit pouvoir savoir, à chaque endroit,
% ce qui était sur le feuillet et ce qui a été ajouté. Les six macros qui
% suivent sont donc l'appareil critique, et non de la décoration.
%
% Les mêmes commandes sont reconnues par `scripts/render.mjs`, qui produit la
% vue de lecture du site. Une macro ajoutée ici doit l'être aussi là-bas,
% faute de quoi le rendu échouera bruyamment — ce qui est voulu : un rendu
% silencieusement faux est pire qu'un rendu absent.

\ProvidesPackage{grothendieck}[2026/08/08 Transcriptions du fonds Grothendieck]

\usepackage{amsmath,amssymb,amsthm}
\usepackage{mathrsfs}
\usepackage{stmaryrd}
% Les diagrammes commutatifs sont partout dans ce fonds : c'est la langue
% même de la géométrie algébrique de Grothendieck, et une transcription qui
% les rendrait en prose ne transcrirait rien.
\usepackage{tikz-cd}
% Français par défaut — c'est la langue des feuillets — mais l'anglais est
% chargé aussi : la traduction et le résumé sont des documents anglais, et la
% typographie française (espaces avant les deux-points) y serait une faute.
% Ces documents basculent par \selectlanguage{english} en tête de corps.
\usepackage[english,french]{babel}
\usepackage{fontspec}
\usepackage{geometry}
\geometry{a4paper,margin=25mm,marginparwidth=28mm}
\usepackage{marginnote}
\usepackage{xcolor}
% ulem, not soul. `soul`'s \st and \ul are built on a character-by-character
% scan that XeLaTeX's font machinery breaks: the document fails with an
% undefined control sequence rather than degrading. ulem does the same two jobs
% and is Unicode-engine safe. `normalem` stops it hijacking \emph, which the
% transcriptions use for Grothendieck's own emphasis.
\usepackage[normalem]{ulem}
\usepackage{hyperref}
\hypersetup{colorlinks=true,linkcolor=black,urlcolor=black,pdfborder={0 0 0}}

% --- Métadonnées du lot -----------------------------------------------------
% Reprises telles quelles par le script de rendu, qui les affiche en tête de
% la vue de lecture. Elles fixent ce que le fichier prétend couvrir : sans
% elles, une transcription flotte sans point d'attache dans un fonds de seize
% mille feuillets.

\newcommand{\folder}[1]{\gdef\grothFolder{#1}}
\newcommand{\batch}[1]{\gdef\grothBatch{#1}}
\newcommand{\leaves}[2]{\gdef\grothFirst{#1}\gdef\grothLast{#2}}
\newcommand{\foldertitle}[1]{\gdef\grothFoldertitle{#1}}
\gdef\grothFolder{?}\gdef\grothBatch{?}\gdef\grothFirst{?}\gdef\grothLast{?}\gdef\grothFoldertitle{}

% --- Le résumé --------------------------------------------------------------
%
% Il ouvre la lecture modernisée, sous le titre « Résumé », et il est composé
% en petit. Ce n'est pas un ornement : c'est le seul passage du document qui
% ne soit pas des mathématiques, et un lecteur doit pouvoir le reconnaître
% avant de l'avoir lu — puis le sauter s'il connaît déjà le sujet, ou s'y
% arrêter s'il ne le connaît pas. Le filet qui le suit marque la reprise du
% corps.

\newenvironment{resume}
  {\small\setlength{\parskip}{0.45em}}
  {\par\medskip\hrule\medskip}

% --- Mots-clés ---------------------------------------------------------------
%
% En anglais, en clôture du résumé de la lecture modernisée : le vocabulaire
% moderne sous lequel chercher ce que les feuillets construisent. Le manifeste
% du site (`npm run manifest`) les extrait pour étiqueter la cote entière —
% cette ligne est la source unique des tags, il n'y a pas de fichier à côté.
\newcommand{\keywords}[1]{\par\smallskip\noindent
  {\footnotesize\textsc{Keywords} --- \textit{#1}}\par}

% --- L'appareil critique ----------------------------------------------------

% Le numéro de feuillet, en marge. C'est l'ancre qui permet de régler un
% désaccord sur une formule en regardant *un* feuillet plutôt qu'une chemise
% de six cents. Le site s'en sert aussi pour tourner les pages du fac-similé
% quand on fait défiler la transcription.
\newcommand{\leaf}[1]{%
  \marginnote{\footnotesize\color{gray!70}\textsf{#1}}%
  \label{leaf:#1}\ignorespaces}

% Illisible : on ne devine pas. Un mot inventé qui se lit comme les autres est
% le pire résultat possible.
\newcommand{\ill}{\textcolor{red!60!black}{[\,\ldots\,]}}

% Lecture douteuse : le mot est donné, et signalé comme incertain.
\newcommand{\uncertain}[1]{\uline{#1}}

% Ajout de l'éditeur — une lettre manquante, un mot sous-entendu.
\newcommand{\add}[1]{\textcolor{blue!50!black}{[#1]}}

% Biffé par Grothendieck lui-même. Ce qu'il a rayé fait partie du document :
% une rature dit souvent où la pensée a changé de direction.
\newcommand{\struck}[1]{\sout{#1}}

% Note du transcripteur, sur l'état matériel du feuillet.
\newcommand{\note}[1]{\par\smallskip{\small\color{gray!60!black}\textit{[#1]}}\par\smallskip}

% Note marginale de Grothendieck, distincte du corps du texte.
\newcommand{\marginal}[1]{\par\smallskip\noindent\hspace{1em}%
  {\small\color{gray!60!black}#1}\par\smallskip}

% --- Titre ------------------------------------------------------------------

\AtBeginDocument{%
  \begin{center}
    {\large\bfseries Fonds Alexandre Grothendieck --- cote n\textsuperscript{o}~\grothFolder}\\[2pt]
    {\itshape\grothFoldertitle}\\[4pt]
    {\small feuillets \grothFirst--\grothLast\ (lot \grothBatch)}\\[2pt]
    {\footnotesize\color{gray!60!black}Transcription établie d'après le fac-similé numérisé
     par l'Université de Montpellier.\\
     Les lectures incertaines sont soulignées, les illisibles notées [\,\ldots\,],
     les ajouts entre crochets.}
  \end{center}
  \vspace{1em}}

\endinput


\folder{117}
\batch{1}
\pages{1}{20}
\dating{1983}
\watermark{Édition de démonstration}
\foldertitle{Théories homotopiques : notes manuscrites (1983, s.d.).}

\begin{document}

\section{Structure des dérivateurs}
\note{titre de sa main, encadré en tête du feuillet 1}

\page{1}
\marginal{déc. 83} \note{souligné, en haut à droite du feuillet}

\struck{$\underline{\mathrm{Hom}}$}

\emph{Donnée} $\mathcal H(I)$,

\[
  (1) \quad \mathrm{Hom}(x, y) \in \mathrm{Ens}
  \qquad\qquad
  \mathrm{Hom}(x, y) \simeq \pi_0(\mathrm{Homhot}(x, y))
\]
\[
  (2) \quad \mathrm{Homhot}(x, y) \in (\mathrm{Hot})
  \qquad\qquad
  \mathrm{Homhot}(x, y) \simeq \int_I \underline{\mathrm{Homhot}}(x, y)
\]
\note{le premier terme de la seconde formule de droite est surchargé ; lu
$\mathrm{Homhot}$ sans certitude}
\[
  (3) \quad \boxed{\underline{\mathrm{Homhot}}(x, y) \in \mathrm{Hot}(I)}
\]
\note{l'indice de $\mathrm{Hot}$ dans l'encadré est surchargé et ne se lit
pas}

Soit $f : I \to J$, \ill{}, pour $X$ dans $\mathcal H(I)$, $Y$ dans
$\mathcal H(J)$,
\[
  (1') \quad \mathrm{Hom}(f_!(X), Y) \simeq \mathrm{Hom}(X, f^{*}Y)
\]
adjonction
\[
  (2') \quad \mathrm{Homhot}(f_!(X), Y) \simeq \mathrm{Homhot}(X, f^{*}(Y))
\]
\uncertain{et plus complète}, \uncertain{donnant} $(1')$ en passant aux
$\pi_0$ \note{lecture de toute la ligne incertaine}
\[
  (3') \quad \boxed{\underline{\mathrm{Homhot}}(f_!X, Y) \simeq f_{*}\,
  \underline{\mathrm{Homhot}}(X, f^{*}(Y))}
\]
donnant $(2')$ en intégrant $\int_J$.

\marginal{duales} \note{ce mot, écrit en long dans la marge gauche, est
relié par un crochet aux deux groupes de formules}
\[
  (1'') \quad \mathrm{Hom}(f^{*}(Y), X) \simeq \mathrm{Hom}(Y, f_{*}(X))
\]
\[
  (2'') \quad \mathrm{Homhot}(f^{*}(Y), X) \simeq \mathrm{Homhot}(Y, f_{*}(X))
\]
\note{dans $(1'')$ et $(2'')$ le signe $\simeq$ est traversé d'un trait
sinueux qui relie les deux lignes ; au-dessus de $(1'')$, un petit
« $=$ » isolé}
\[
  (3'') \quad \boxed{\underline{\mathrm{Homhot}}(Y, f_{*}(X)) \simeq f_{*}\,
  \underline{\mathrm{Homhot}}(f^{*}Y, X)}
\]
\note{dans $(3'')$ le premier argument est d'abord écrit, puis biffé et
remplacé par $Y$ ; sous l'encadré, un $f_*$ isolé, comme une étiquette}

\page{2}
\note{brouillon sur le même papier de listing, sans phrase ; les formules
sont données dans l'ordre de la page, les croquis décrits}
\[
  \pi_0(x, y) = \mathrm{Hom}_{\mathcal H}(x, y) \qquad
  \pi_1(x, y ; f) = \struck{\ill{}}
\]
\[
  x \xrightarrow{\ f\ } y \qquad
  \mathcal H(\Delta_1) \to \mathcal H(\Delta_0) \times \mathcal H(\Delta_0),
  \quad \varphi \longmapsto
\]
\[
  J \hookrightarrow I \qquad X_J \qquad \mathcal H(J) \qquad
  \mathcal H(I, J) \qquad \mathcal H(I) \to \mathcal H(J)
\]
\[
  X_i \to X_J, \qquad Y_i \to Y_J
\]
\[
  \pi_1(x, y ; f, g)
\]
\note{ce dernier symbole est précédé d'un croquis : deux flèches $f$ et $g$
de $x$ vers $y$, une double flèche verticale entre elles ; ailleurs sur la
page, plusieurs croquis de 2-flèches entre deux flèches $x \rightrightarrows
y$, et un cycle de flèches alternées entre $x$ et $y$}

\[
  \mathcal H(\Delta_1) \qquad \longrightarrow\longrightarrow\longrightarrow
  \qquad M \ W \qquad \mathrm{Hot}
\]

\struck{$\pi_W(x, y) =$}

\struck{$\underline{\mathrm{Hom}}$} \quad $\mathrm{Cl}\,(M, W)$ \ill{}
\note{au pied de la feuille, rognée, une ligne où se lisent seulement
$M$ et $(\mathrm{Cat})$}

\page{3}
\begin{enumerate}
  \item[a)] Modèles --- leur importance, leur indétermination.
  \item[b)] Suites exactes longues (évt.\ de deux types) associées à une
  flèche --- liées à la construction des cylindre et cocylindre d'un
  objet, modèle, et \struck{\ill{}} des cylindres et cocylindre d'une
  application entre modèles.
  \item[c)] Notion de résolution d'un modèle (gauche ou droite) formalisme
  des foncteurs dérivés (quand il y a deux th.\ homotopiques \add{en
  présence}) \note{« en présence » est écrit sous la ligne ; la lettre c)
  est cerclée}
  \item[d)] Notions homotopiques sur modèles (structures homotopiques :
  $W$, \uncertain{cofibr.}, homotopismes…) « Résolutions définies à
  homotopie près ».
  \marginal{mais ce sont des notions intermédiaires \uncertain{pour} la
  construction des cat.\ dérivées}
  \item[e)] Homotopies entre applications, homotopies entre homotopies
  etc. On a sans doute \struck{\ill{}}, pour deux \struck{modèles}
  $X, Y$ un objet $\mathrm{Hom}(X, Y)$ dans $\mathrm{Hot}$ \add{ou
  « simplicial »} \struck{\ill{}} $=$ un complexe de type convenable,
  explicitant \ill{} leur donnant naissance \ill{} dans $\mathcal H$ aux
  $\mathrm{Hom}$ extérieurs à valeurs en \ill{} $\mathrm{Hot}$ ---, ou
  dans une théorie plus riche $\mathrm{Hot}$ --- (dans Quillen, il y a
  des $\mathrm{Hom}$ simpliciaux)
  \note{« ou “simplicial” » est interlinéaire, sous
  $\mathrm{Hom}(X, Y)$ ; la phrase qui suit « convenable » ne se lit que
  par morceaux}
  \item[] \struck{f) Cf.\ plus bas}
  \item[\uncertain{f)}] Opérations de \ill{}, soit \uncertain{grossières}
  sur modèles, soit délicates sur $\mathcal H$.
  \item[\uncertain{g)}] Invariants $\pi_i$ ou $H_i$, à valeurs dans une
  catégorie (souvent abélienne) associée à la théorie.
\end{enumerate}
\note{les lettres de ces deux derniers items sont surchargées à l'encre
noire au point de ne plus se lire ; elles sont restituées d'après la note
marginale, qui les nomme}
\marginal{f) et g) sont des \uncertain{aspects} particulièrement peu
compris}

\page{4}
\begin{enumerate}
  \item[\uncertain{h)}] Possibilité de définir \struck{des} invariants
  $K(\mathcal H)$ via « triangles exacts », voir $K_i$
  \uncertain{supérieurs} d'une théorie homotopique. (Cas additif
  seulement ?!?)
  \item[\uncertain{i)}] Souvent il y a \add{des} $\underline{\mathrm{Hom}}$
  internes dans $\mathcal H$, avec \struck{des} \add{une} formules du type
  \[
    \mathrm{Hom}(\xi, \eta) = \Gamma\, \underline{\mathrm{Hom}}(\xi, \eta)
  \]
  \[
    \Gamma \text{ ou } \pi_0 : \mathcal H \to \mathrm{Ens}
  \]
  foncteur important (souvent $\mathrm{Hom}(1, -)$, avec $1$ un obj.\
  final…) Souvent aussi, structure multiplicative (soit produit
  cartésien, soit un produit $\otimes$) avec une formule
  \[
    \mathrm{Hom}(x \otimes y, z) \simeq \mathrm{Hom}(x,
    \underline{\mathrm{Hom}}(y, z))
  \]
\end{enumerate}
\note{les deux lettres d'item sont surchargées (un $h$ sur un $i$, un $i$
sur un $h$, semble-t-il) ; une longue flèche, dans la marge gauche, part du
second item vers le haut de la page, vers le premier ; « (des) » est
interlinéaire, au-dessus de « Hom »}
\marginal{Structure multiplicative et structure $\mathrm{Hom}$}
\marginal{(24.11.83) relations entre $\underline{\mathrm{Hom}}$ internes
et les $\mathrm{Homhot}(x, y) \in \mathrm{Ob}\,\mathrm{Hot}$ dans e) ??}

Sauf peut-être les structures f) et g)), il semble que tout puisse se
décrire en termes de $(M, W \subset \mathrm{Fl}(M))$.

\emph{Dégager les bons axiomes}

\section{Théorie (de variance) homotopique $\mathcal H$}
\note{première ligne du feuillet 5, de sa main ; c'est le début d'une suite
qu'il pagine lui-même de 1 à 9, au milieu du bord supérieur}

\page{5}
\note{p.~1 de l'auteur}
\marginal{\emph{(1) Variance}}
Elle associe à toute \struck{\ill{}} \add{petite} catégorie $A$
\struck{$\in \mathrm{Ob}\,(\mathrm{Cat})$} \struck{(ou non petite)} une
catégorie (non néc.\ petite \struck{\ill{} si $A$ l'est})
\[
  (1) \qquad \boxed{A \longmapsto \mathcal H(A)}
\]
\struck{et de plus que} \uncertain{qui} \struck{\ill{}} 2-contravariante
en $A$ (de façon \emph{stricte})
\[
  (2) \qquad
  \begin{array}{c} f : A \to B \\ \text{foncteur entre $U$-catégories}
  \end{array}
  \quad \text{implique} \quad
  \mathcal H(B) \xrightarrow{\ f^{*}_{\mathcal H}\ } \mathcal H(A)
\]
(\ill{}). On pose
\[
  (3) \qquad \mathcal H(\Delta_0) = \mathrm{Hot}_{\mathcal H}
\]
\ill{} \struck{\ill{}} on a donc pour $A$ un foncteur
\[
  (4) \qquad \mathcal H(A) \xrightarrow{\ \alpha_A\ }
  \underline{\mathrm{Hom}}(A, \struck{\mathrm{Hot}}\ \mathrm{Hot}_{\mathcal H}),
  \qquad
  \mathcal X \longmapsto \big(a \longmapsto (i_a)^{*}_{\mathcal H}(\mathcal X)\big)
\]
où
\[
  i_a : \Delta_0 \longrightarrow A \quad \text{foncteur de valeur } a
\]
(\emph{NB} Pour tt $B$, on a
\[
  \underline{\mathrm{Hom}}(B, A) \longrightarrow
  \underline{\mathrm{Hom}}(\mathcal H(A), \mathcal H(B))
\]
d'où
\[
  \mathcal H(A) \longrightarrow \underline{\mathrm{Hom}}(\underline{\mathrm{Hom}}(B, A),
  \mathcal H(B))
\]
si $B = \Delta_0$, on trouve (4).)

Pour $f : A \to B$ morphisme de catégories, on a donc

\page{7}
\note{p.~2 de l'auteur}

\begin{tikzcd}
  \mathcal H(B) \arrow[r, "\mathcal H(f)"] \arrow[d, "\alpha_B"'] & \mathcal H(A) \arrow[d, "\alpha_A"] \\
  \underline{\mathrm{Hom}}(B, \mathrm{Hot}_{\mathcal H}) \arrow[r, "f^{*}_{\mathrm{Hot}_{\mathcal H}}"'] & \underline{\mathrm{Hom}}(A, \mathrm{Hot}_{\mathcal H})
\end{tikzcd}

\note{carré numéroté (5) en marge ; l'étiquette de la flèche du bas est
écrite sous elle, avec une parenthèse fermante}

\emph{(Hot 1)} \quad si $A = \coprod_{i \in I} A_i$, alors
$\mathcal H(A) \xrightarrow{\ \sim\ } \prod_{i \in I} \mathcal H(A_i)$
\emph{iso}

\textbf{(2)} \quad $\int$ (intégration).
\note{le 2 est cerclé, comme « (Hot 1) »}

Pour \struck{toute} \add{toute flèche $f : A \to B$} \struck{$A$ petite,
i.e.\ $A \in \mathrm{Ob}\,(\mathrm{Cat})$}, on donne
\[
  (6) \qquad \boxed{\mathcal H(A) \xrightarrow{\ \underline{L}_A\ }
  \mathcal H(A^{\wedge})}
\]
\uncertain{rendant} commutatif \ill{}
\[
  (7) \qquad \mathcal H(A) \xrightarrow{\ \underline{L}_A\ }
  \mathcal H(A^{\wedge}) \xrightarrow{\ (\varepsilon_A)^{**}_{\mathcal H}\ }
  \mathcal H(A),
  \qquad \text{composé} = \mathrm{id}_{\mathcal H(A)}
\]
où
\[
  \varepsilon_A : A \hookrightarrow A^{\wedge}
\]
\uncertain{est} l'inclusion canonique.
\marginal{\uncertain{(Hot 2)}} \note{souligné, en marge de (7) ; au-dessus
de $(\varepsilon_A)^{**}_{\mathcal H}$, un mot biffé illisible}

\emph{Hot 3} \quad $\forall f : A \to B$ dans $(\mathrm{Cat})$,
$\exists f_! : \mathcal H(A) \to \mathcal H(B)$, \ill{} \uncertain{à
isomorphisme près}, tel qu'on ait commutativité

\begin{tikzcd}
  \mathcal H(A) \arrow[r, "\underline{L}_A"] \arrow[d, "f_!"'] & \mathcal H(\widehat{A}) \arrow[d, "(f^{*}_{\wedge})^{*}_{\mathcal H} \overset{\mathrm{def}}{=} \hat f_!"] \\
  \mathcal H(B) \arrow[r, "\underline{L}_B"'] & \mathcal H(B^{\wedge})
\end{tikzcd}

\note{carré numéroté (8) en marge. Tout le passage de (6) à (8) est barré de
longs traits obliques : c'est une première version de l'intégration,
abandonnée, que la page 9 remplace par (Hot 2)}

\page{9}
\note{p.~3 de l'auteur}
On suppose

\emph{(Hot 2)} \quad $\forall f : A \to B$ dans $\mathrm{Cat}$, le foncteur
\struck{$f^{*} : \mathrm{Hot}(A)$}
\[
  f^{*}_{\mathcal H} : \mathcal H(B) \to \mathcal H(A)
\]
a un \emph{adjoint à gauche}
\[
  f^{\mathcal H}_{!} : \mathcal H(B) \to \mathcal H(A)
\]
\note{sic : le but et la source de $f^{\mathcal H}_!$ sont écrits comme ceux
de $f^{*}_{\mathcal H}$}

$\mathcal H(A)$ \ill{} \emph{\ill{}} \ill{} \note{une ligne et demie qui
ne se lit pas}

\textbf{A.} Il faudra faire attention, s'il ne faudrait pas
\uncertain{nous} borner ici aux \struck{$\ast$} $f$
\emph{\uncertain{cofibrantes}}.
\note{la fin du mot souligné est surchargée}

Si $B = \Delta_0$, on trouve
\[
  \mathcal H(A) \longrightarrow \mathrm{Hot}_{\mathcal H}
\]
noté $\int_A$, et appeler foncteur d'intégration,
\marginal{i.e.\ $\pi_A : A \to \Delta_0$, le \uncertain{prolongement}
$(\pi_A)_!$}
jouant le rôle de un $\varinjlim_A$.

On a \ill{}, par \uncertain{transitivité},
\[
  \int_A \mathcal X = \int_B f_!(\mathcal X)
\]
\uncertain{pour} $f : A \to B$.
\marginal{\emph{NB} L'existence des $f_!$ implique l'existence des sommes
dans les $\mathcal H(A)$, \uncertain{et les $f^{*}$ bien sûr y
commutent}}

\textbf{(3)} \quad Opérations $\prod_A$ de cointégration, \ill{} des
$f^{\mathcal H}_{*}$ adjoints à droite des $f^{*}_{\mathcal H}$.
\marginal{(Hot 3)} \note{le 3 et « Hot 3 » sont cerclés}

\page{11}
\note{p.~4 de l'auteur}
On déduit des foncteurs
\[
  \prod_A = (\pi_A)_{*} : \mathcal H(A) \longrightarrow \mathrm{Hot}_{\mathcal H}
\]
\note{l'égalité est écrite verticalement, $(\pi_A)_*$ sous $\prod_A$}
\[
  \struck{\ast}\ \prod_A \mathcal X \simeq \prod_B f_{*}(\mathcal X).
\]
\marginal{L'existence des $f^{\mathcal H}_{*}$ implique celle des produits
dans les $\mathcal H(A)$, les $f^{*}$ y commutent.}

\emph{Dualité} \quad Posons $\mathcal H^{\circ}(A) = \mathcal H(A)^{\circ}$ ;
on trouve \ill{}
les
\[
  f^{*}_{\mathcal H^{\circ}} = (f^{*}_{\mathcal H})^{\circ} :
  \mathcal H^{\circ}(B) \to \mathcal H^{\circ}(A),
\]
et
\[
  \begin{cases}
    f^{\mathcal H^{\circ}}_{!} = (f^{\mathcal H}_{*})^{\circ} \\
    f^{\mathcal H^{\circ}}_{*} = (f^{\mathcal H}_{!})^{\circ}
  \end{cases}
\]
Les opérations $\int$ et $\prod$ s'échangent.

\emph{Cas additif} \quad les $\mathcal H(A)$ \struck{\ill{}} cat.\
additives (resp.\ $k$-linéaires, $k$ anneau fixé), $f^{*}_{\mathcal H}$
additifs, alors les $f^{\mathcal H}_{!}$, $f^{\mathcal H}_{*}$ additifs.

\page{13}
\note{p.~5 de l'auteur}
\emph{Ex 1} \quad $M$ catégorie, avec $\varinjlim$ et $\varprojlim$
quelconques. On pose
\[
  \mathcal H(A) = \underline{\mathrm{Hom}}(A, M)
\]
Alors on sait que les foncteurs
\[
  f^{*}_{M} : \underbrace{\underline{\mathrm{Hom}}(B, M)}_{\mathcal H(B)}
  \longrightarrow \underbrace{\underline{\mathrm{Hom}}(A, M)}_{\mathcal H(A)}
\]
associés aux $f : A \to B$ dans $(\mathrm{Cat})$, admettent des adjoints : à
g.\ $f^{M}_{!}$ et à dr.\ $f^{M}_{*}$.
\note{les égalités $= \mathcal H(B)$ et $= \mathcal H(A)$ sont écrites
verticalement sous les deux $\underline{\mathrm{Hom}}$}

Ici $\mathrm{Hot}_{\mathcal H} = M$ \quad condition triviale. Ici
\[
  \begin{cases}
    \int_A = \varinjlim_A \\
    \prod_A = \varprojlim_A
  \end{cases}
\]
--- \uncertain{c'est bien connu}.

\emph{Ex 2} \quad Soit $(M, W)$, avec $W \subset \mathrm{Fl}(M)$. On
définit
\[
  W_A \subset \mathrm{Fl}\,\underline{\mathrm{Hom}}(A, M)
\]
\struck{\ill{}} \ill{} fibre par fibre, et on pose
\[
  \mathcal H(A) = W_A^{-1}\, \underline{\mathrm{Hom}}(A, M)
\]
\uncertain{l'existence des} $f^{*}_{\mathcal H}$ est évidente.
\marginal{(Hot 1) \uncertain{est satisfait}, puisque \ill{}
$\mathcal H(\coprod A_i) \simeq \prod_{i \in I} \mathcal H(A_i)$
\ill{}}

\page{15}
\note{p.~6 de l'auteur}
On s'intéresse aux cas où les $f^{*}_{\mathcal H}$ ont des adjoints : à
g.\ $f^{\mathcal H}_{!}$ et à droite $f^{\mathcal H}_{*}$, une condition sur
$W$.

\emph{Remarque} \quad Soit
\[
  \underset{(W)}{M} \xrightarrow{\ f\ } \underset{(W')}{M'}
\]
tel que $f(W) \subset W'$. Alors
\[
  \mathcal H_W(A) \Longrightarrow \mathcal H_{W'}(A)
\]
homomorphisme \uncertain{fonctoriel} en $A$, (i.e.\ \uncertain{compatible}
aux $f^{*}_{W}$, $f^{*}_{W'}$), on ne sait \uncertain{nullement} pas a priori
qu'il « commute » \ill{} aux $f_!$, $f_*$. Par exemple si $W, W'$ sont
réduites aux iso, dans \ill{} cette commutation \struck{signifie}
\ill{} \add{(resp.\ \ill{})} que $f$ commute aux
$\varinjlim$ (resp.\ $\varprojlim$). Mais bien sûr si
\[
  \mathcal H_W(A) \to \mathcal H_{W'}(A)
\]
est une équivalence
\note{la phrase se poursuit en tête de la page 17 ; l'argument de
$\mathcal H_{W'}$ porte un exposant illisible}

\page{17}
\note{p.~7 de l'auteur}
pour tt $A$, ça sera OK.

Voici un cas où il en sera ainsi : Supposons que $f$ admette un adjoint à
droite (disons) $g$ satisfaisant aux conditions
\[
  \text{(i)} \quad f^{-1}(W') = W
\]
\[
  \text{(ii)} \quad fg(X') \to X' \in W' \qquad \forall\, X' \in \mathrm{Ob}\, M'
\]
Alors
\[
  \underline{\mathrm{Hom}}(A, M)
  \underset{g^{M}}{\overset{f^{M}}{\rightleftarrows}}
  \underline{\mathrm{Hom}}(A, M')
\]
sont adjoints l'un de l'autre, et \struck{les} \ill{} satisfont aux
\add{mêmes} conditions pour $W_A, W'_A$. \struck{Donc}
\note{« mêmes » est écrit « $=$s » ; dans les deux étiquettes $f^M$, $g^M$,
une étoile est biffée}
\[
  \mathcal H_W(A)
  \underset{g^{A}_{W}}{\overset{f^{A}_{W}}{\rightleftarrows}}
  \mathcal H_{W'}(A)
\]
\struck{\ill{}}
\note{l'indice de $\mathcal H_W$ et du second terme sont surchargés ; la
ligne qui suit le diagramme est biffée de plusieurs traits et ne se lit
pas}

Les théories $\mathcal H$ qui peuvent se décrire par les $(M, W)$ sont
dites \emph{spéciales}. Cela implique des conditions draconiennes, sur
$(M, W)$

\page{19}
\note{p.~8 de l'auteur}
On peut se \struck{\ill{}} \add{proposer d'obtenir} une vue d'ensemble sur
les $(M, W)$ qui donnent naissance à une théorie homotopique donnée
$\mathcal H$. \struck{S'il y a \ill{}} Pour deux tels, $(M_1, W_1)$,
$(M_2, W_2)$, peut-on trouver une troisième $(M, W)$, et \struck{une
\ill{}}
\[
  \begin{array}{c}
    (M_1, W_1) \longrightarrow (M, W) \\
    (M_2, W_2) \nearrow
  \end{array}
\]
donnant naissance à l'iso.\ \struck{\ill{}} donné entre $\mathcal H_{W_1}$
et $\mathcal H_{W_2}$ ?
\note{les deux flèches vers $(M, W)$ sont dessinées l'une sous l'autre, la
seconde oblique ; « proposer d'obtenir » et « $(M_1,W_1), (M_2,W_2)$, pour
tels » sont en interligne}

\emph{Ex 3} \quad (Cas particulier)

$M = (\mathrm{Cat})$, $W$ un localisateur fondamental.

J'ai prouvé l'existence des $f_!$, mais non encore celle des
$f_*$.

Il est possible qu'il faille pour les \ill{} faire des hypothèses de
finitude sur les $A, B$…

\page{20}
\note{p.~9 de l'auteur}
\emph{Ex 4} \quad Autres cas particuliers. Soit
$M =$ complexes d'une catégorie abélienne, et $W$ les quasi-isom.
\[
  \underline{\mathrm{Hom}}(A, M) = \text{complexes de }
  \underline{\mathrm{Hom}}(A, \mathcal A)
\]
\[
  W_A = \text{quasi-isom.\ d'icelle.}
\]
$\mathcal H(A)$ catégorie dérivée de $\underline{\mathrm{Hom}}(A, \mathcal A)$.
\note{cette dernière ligne est interlinéaire, reliée par un crochet à
$W_A$}

Si p.ex.\ $\mathcal A =$ faisceaux de $k$-modules sur topos
\struck{\ill{}} $\mathcal T$, alors
\[
  \underline{\mathrm{Hom}}(A, \mathcal A) = \text{faisceaux de $k$-modules
  sur topos } A^{\circ\wedge} \times_{\mathrm{top}} \mathcal T
\]
et
\[
  f^{*} : \underline{\mathrm{Hom}}(B, M) \to \underline{\mathrm{Hom}}(A, M)
\]
est le foncteur \uncertain{image inverse} associé :

\begin{tikzcd}
  A^{\circ\wedge} \times \mathcal T \arrow[d, "f^{\circ}_{\mathrm{top}} \times \mathrm{id}_{\mathcal T}"] \\
  B^{\circ\wedge} \times \mathcal T
\end{tikzcd}

\note{à droite du diagramme, « \uncertain{l'induit est} $(\mathrm{L}
f^{\circ *})$ », lecture incertaine}

Le foncteur \add{induit sur les catégories dérivées} \struck{a un adjoint
à droite} a un adjoint à \emph{droite} $\mathrm{R}f_{*}$, \struck{\ill{}}
généralement pas d'adjoint à gauche \ill{} \uncertain{sauf} si
$\mathcal T =$ topos ponctuel i.e.\ $\mathcal A = (k\text{-Modules})$. Si
pourtant, on \uncertain{se place} au niveau des \struck{topos} faisceaux
d'ensembles, $(f^{\circ}_{\mathrm{top}})_!$ existe et est exact à gauche,
donc au niveau des \add{complexes de} faisceaux abéliens il
\struck{\ill{}} transforme quasi-isom.\ en q-isom., donc passe aux
catégories dérivées en $\mathrm{L}f_!$.
\note{« induit sur les catégories dérivées » est en interligne, au-dessus de
« à droite » ; « complexes de » aussi}

\emph{Propriétés formelles} \quad \struck{\ill{}}
\add{\uncertain{les commutations des foncteurs produits}}

\emph{(Hot 4)}

\begin{tikzcd}
  A \arrow[r, "g'"] \arrow[d, "f"'] & A' \arrow[d, "f'"] \\
  B \arrow[r, "g"'] & B'
\end{tikzcd}

\[
  g^{*} f'_{!} \xleftarrow{\ \sim\ } f_{!}\, g'^{*}
  \quad \text{si $f$ cofibrant [ou $g$ fibrant ??]}
\]
\[
  g^{*} f'_{*} \xrightarrow{\ \sim\ } f_{*}\, g'^{*}
  \quad \text{si $f$ fibrant [ou $g$ cofibrant ??]}
\]
\note{au centre du carré, « cart. » ; au-dessus de la première formule, deux lignes d'interligne en partie
lisibles, où se lit « de $g^*$ » ; la flèche de la première est un
$\simeq$ prolongé d'une pointe vers la gauche, lue $\xleftarrow{\sim}$ sans
certitude}

[Les formules \uncertain{sont} « transposées » \uncertain{l'une de l'autre}
(par passage aux adjoints à droite, resp.\ à g.)]

\emph{(Hot 5)} \quad Si $A \xrightarrow{\ f\ } B$ pl.\ fid., alors
\[
  f^{*}_{\mathcal H} : \mathcal H(B) \to \mathcal H(A) \quad \text{foncteur
  localisation}
\]
(i.e.\
\[
  f^{\mathcal H}_{*},\ f^{\mathcal H}_{!} : \mathcal H(A) \to \mathcal H(B)
  \quad \text{sont pl.\ fid.}
\]
et
\[
  f^{*} f_{!} \simeq \mathrm{id}, \qquad f^{*} f_{*} \simeq \mathrm{id}\,)
\]

\emph{(Hot 6)} \quad Si $f : A \to B$ localisation, alors
\[
  f^{*}_{\mathcal H} : \mathcal H(B) \to \mathcal H(A) \quad \text{pl.\ fid.,
  i.e.}
\]
\[
  f^{\mathcal H}_{!},\ f^{\mathcal H}_{*} \quad \text{foncteurs loc.}
  \quad \struck{\text{iso}} \quad \text{et} \quad
  f_! f^{*} \simeq \mathrm{id} \simeq f_{*} f^{*}
\]
\note{la dernière formule est écrite au pied du feuillet, serrée contre le
bord ; lue sans certitude}

\end{document}
